\documentclass[12pt,letterpaper]{article}
\usepackage{graphicx,textcomp}
\usepackage{natbib}
\usepackage{setspace}
\usepackage{fullpage}
\usepackage{color}
\usepackage[reqno]{amsmath}
\usepackage{amsthm}
\usepackage{fancyvrb}
\usepackage{amssymb,enumerate}
\usepackage[all]{xy}
\usepackage{endnotes}
\usepackage{lscape}
\newtheorem{com}{Comment}
\usepackage{float}
\usepackage{hyperref}
\newtheorem{lem} {Lemma}
\newtheorem{prop}{Proposition}
\newtheorem{thm}{Theorem}
\newtheorem{defn}{Definition}
\newtheorem{cor}{Corollary}
\newtheorem{obs}{Observation}
\usepackage[compact]{titlesec}
\usepackage{dcolumn}
\usepackage{tikz}
\usetikzlibrary{arrows}
\usepackage{multirow}
\usepackage{xcolor}
\newcolumntype{.}{D{.}{.}{-1}}
\newcolumntype{d}[1]{D{.}{.}{#1}}
\definecolor{light-gray}{gray}{0.65}
\usepackage{url}
\usepackage{listings}
\usepackage{color}

\definecolor{codegreen}{rgb}{0,0.6,0}
\definecolor{codegray}{rgb}{0.5,0.5,0.5}
\definecolor{codepurple}{rgb}{0.58,0,0.82}
\definecolor{backcolour}{rgb}{0.95,0.95,0.92}

\lstdefinestyle{mystyle}{
	backgroundcolor=\color{backcolour},   
	commentstyle=\color{codegreen},
	keywordstyle=\color{magenta},
	numberstyle=\tiny\color{codegray},
	stringstyle=\color{codepurple},
	basicstyle=\footnotesize,
	breakatwhitespace=false,         
	breaklines=true,                 
	captionpos=b,                    
	keepspaces=true,                 
	numbers=left,                    
	numbersep=5pt,                  
	showspaces=false,                
	showstringspaces=false,
	showtabs=false,                  
	tabsize=2
}
\lstset{style=mystyle}
\newcommand{\Sref}[1]{Section~\ref{#1}}
\newtheorem{hyp}{Hypothesis}


\title{Problem Set 4}
\date{Due: November 26, 2026}
\author{Applied Stats/Quant Methods 1}


\begin{document}
	\maketitle
	\section*{Instructions}
	\begin{itemize}
		\item Please show your work! You may lose points by simply writing in the answer. If the problem requires you to execute commands in \texttt{R}, please include the code you used to get your answers. Please also include the \texttt{.R} file that contains your code. If you are not sure if work needs to be shown for a particular problem, please ask.
		\item Your homework should be submitted electronically on GitHub.
		\item This problem set is due before 23:59 on Thursday November 26, 2026. No late assignments will be accepted.
	\end{itemize}



	\vspace{.5cm}
\section*{Question 1: Economics}
\vspace{.25cm}
\noindent 	
In this question, use the \texttt{prestige} dataset in the \texttt{car} library. First, run the following commands:

\begin{verbatim}
install.packages(car)
library(car)
data(Prestige)
help(Prestige)
\end{verbatim} 


\noindent We would like to study whether individuals with higher levels of income have more prestigious jobs. Moreover, we would like to study whether professionals have more prestigious jobs than blue and white collar workers.

\newpage
\begin{enumerate}
	
	\item [(a)]
	Create a new variable \texttt{professional} by recoding the variable \texttt{type} so that professionals are coded as $1$, and blue and white collar workers are coded as $0$ (Hint: \texttt{ifelse}). \\
	
	\lstinputlisting[language=R, firstline=53, lastline=53]{PS04_NF.R}
	
	We check that the re-coding was correctly done by opening the data.
	
	
	\vspace{1cm}
	
	
	\item [(b)]
	Run a linear model with \texttt{prestige} as an outcome and \texttt{income}, \texttt{professional}, and the interaction of the two as predictors (Note: this is a continuous $\times$ dummy interaction.) \\
	We run the regression using the built-in R commands:
	\lstinputlisting[language=R, firstline=61, lastline=61]{PS04_NF.R}
	
	\vspace{1cm}
	\item [(c)]
	Write the prediction equation based on the result.
	We look at the results of the regression to write the equation: \\
		\begin{verbatim}
		> reg1
		
		Call:
		lm(formula = prestige ~ income + professional + income:professional, 
		data = Prestige)
		
		Coefficients:
		(Intercept)               income         professional  income:professional  
		21.142259             0.003171            37.781280            -0.002326
	\end{verbatim}
	The equation is therefore: \\
	Prestige = 21.142259 + income*0.003171 + professional*37.781280 - 0.002326*income*professional \\
	\vspace{.5cm} \\
	or, considering the fact that "professional" is a dummy variable: \\
	\vspace{.5cm} \\
	Prestige for professionals = 58.923539  + income(0.003171 - 0.002326) \\
	and\\
	Prestige for white and blue collar = 21.142259 + income*0.003171
	
	
	
\newpage
	\item [(d)]
	Interpret the coefficient for \texttt{income}. \\
	Holding the job category constant, each additional Canadian dollar (1971) of income a person receives is associated with an increase of 0.003171 units in their prestige score.
	
	\vspace{1cm}	
	\item [(e)]
	Interpret the coefficient for \texttt{professional}. \\
	Holding the income constant, a person holding a professional type of job is expected to score 37.781280 units higher than a person having a white or blue collar job, on average.\\
	The interactive effect of income and profession type decreases this prestige score of 0.002326 on average.
	
	\vspace{1cm}
	\item [(f)]
	What is the effect of a \$1,000 increase in income on prestige score for professional occupations? In other words, we are interested in the marginal effect of income when the variable \texttt{professional} takes the value of $1$. Calculate the change in $\hat{y}$ associated with a \$1,000 increase in income based on your answer for (c). \\
	
	We start by calculating the effect of a one-unit increase in income on prestige score for professional occupations:\\
	
	\lstinputlisting[language=R, firstline=95, lastline=96]{PS04_NF.R} 
	\newpage
	We get the following result:
	
		\begin{verbatim}
		> qF_oneUnit
		[1] 0.000845
	\end{verbatim}
	
	The marginal effect of income on prestige score for professional occupations is 0.000845. \\
	Now we multiply this by a thousand:
	\lstinputlisting[language=R, firstline=98, lastline=99]{PS04_NF.R}
	
	We get the following result:
	
	\begin{verbatim}
	> qF_thou
	[1] 0.845
	\end{verbatim}
	
	The effect of a \$1,000 increase in income on prestige score for professional occupations is 0.845: each increase of \$1,000 in a professional's income is associated with an increase of 0.845 increase in their prestige score.
	
	\vspace{1cm}
	
	
	\item [(g)]
	What is the effect of changing one's occupations from non-professional to professional when her income is \$6,000? We are interested in the marginal effect of professional jobs when the variable \texttt{income} takes the value of $6,000$. Calculate the change in $\hat{y}$ based on your answer for (c). \\
	
	We calculate the prestige score for a non-professional earning \$6,000:
	\lstinputlisting[language=R, firstline=111, lastline=112]{PS04_NF.R} 
	
	We get the following result:
	
	\begin{verbatim}
	> qG_nonpro
	[1] 40.16826
	\end{verbatim}
	A non-professional earning \$6,000 has a prestige score of around 40.16 units.
	
	We now calculate the prestige score for a professional earning \$6,000:
	\lstinputlisting[language=R, firstline=115, lastline=116]{PS04_NF.R} 
	
	We get the following result:
	
	\begin{verbatim}
		> qG_pro
		[1] 63.99354
	\end{verbatim}
	A professional earning \$6,000 has a prestige score of around 63.99 units.
	
	We now calculate the difference:	
	\lstinputlisting[language=R, firstline=119, lastline=120]{PS04_NF.R} 
	
	We get the following result:
	
	\begin{verbatim}
	> qG
	[1] 23.82528
	\end{verbatim}
	A change in one's occupation from non-professional to professional when her income is \$6,000 is associated with a 23.82528 units increase in prestige score.
	
	
	
	
\end{enumerate}

\newpage

\section*{Question 2: Political Science}
\vspace{.25cm}
\noindent 	Researchers are interested in learning the effect of all of those yard signs on voting preferences.\footnote{Donald P. Green, Jonathan	S. Krasno, Alexander Coppock, Benjamin D. Farrer,	Brandon Lenoir, Joshua N. Zingher. 2016. ``The effects of lawn signs on vote outcomes: Results from four randomized field experiments.'' Electoral Studies 41: 143-150. } Working with a campaign in Fairfax County, Virginia, 131 precincts were randomly divided into a treatment and control group. In 30 precincts, signs were posted around the precinct that read, ``For Sale: Terry McAuliffe. Don't Sellout Virgina on November 5.'' \\

Below is the result of a regression with two variables and a constant.  The dependent variable is the proportion of the vote that went to McAuliff's opponent Ken Cuccinelli. The first variable indicates whether a precinct was randomly assigned to have the sign against McAuliffe posted. The second variable indicates
a precinct that was adjacent to a precinct in the treatment group (since people in those precincts might be exposed to the signs).  \\

\vspace{.5cm}
\begin{table}[!htbp]
	\centering 
	\textbf{Impact of lawn signs on vote share}\\
	\begin{tabular}{@{\extracolsep{5pt}}lccc} 
		\\[-1.8ex] 
		\hline \\[-1.8ex]
		Precinct assigned lawn signs  (n=30)  & 0.042\\
		& (0.016) \\
		Precinct adjacent to lawn signs (n=76) & 0.042 \\
		&  (0.013) \\
		Constant  & 0.302\\
		& (0.011)
		\\
		\hline \\
	\end{tabular}\\
	\footnotesize{\textit{Notes:} $R^2$=0.094, N=131}
\end{table}

\vspace{.5cm}
\begin{enumerate}
	\item [(a)] Use the results from a linear regression to determine whether having these yard signs in a precinct affects vote share (e.g., conduct a hypothesis test with $\alpha = .05$).\\
	
	Hypotheses:\\
	H0: \( \beta = 0) \)\\
	H1: \( \beta \neq 0) \) \\
	\vspace{.5cm} \\
	We can now calculate our t-statistic, using the beta estimate in the table (0.042) and the standard error (0.016):\\
	
	\lstinputlisting[language=R, firstline=155, lastline=156]{PS04_NF.R} 
	
	We get the following result:
	
	\begin{verbatim}
	> t_stat1
	[1] 2.625
	\end{verbatim}
	
	Then we calculate the degrees of freedom:
	\lstinputlisting[language=R, firstline=158, lastline=159]{PS04_NF.R} 
	
	We get the following result:
	
	\begin{verbatim}
	> df
	[1] 130
	\end{verbatim}
	
	We use these values to conduct a two-tailed t-test:
	\lstinputlisting[language=R, firstline=161, lastline=161]{PS04_NF.R} 
	
	We get the following result:
	
	\begin{verbatim}
		> 2*pt(t_stat1, df, lower.tail=FALSE)
		[1] 0.009703405
	\end{verbatim}
	
	0.001563955 \textless 0.05 \\ 
	We can reject the null hypothesis that our coefficient is 0,
	meaning there is an association between having these yard signs in a precinct and the vote share for Ken Cuccinelli.
	
	
	\newpage		
	\item [(b)]  Use the results to determine whether being
	next to precincts with these yard signs affects vote
	share (e.g., conduct a hypothesis test with $\alpha = .05$). \\
	
	
		Hypotheses:\\
	H0: \( \beta = 0) \)\\
	H1: \( \beta \neq 0) \) \\
	\vspace{.5cm} \\
	We can now calculate our t-statistic, using the beta estimate in the table (0.042) and the standard error (0.013):\\
	
	\lstinputlisting[language=R, firstline=176, lastline=177]{PS04_NF.R} 
	
	We get the following result:
	
	\begin{verbatim}
		> t_stat2
		[1] 3.230769
	\end{verbatim}
	
	Then we calculate the degrees of freedom:
	\lstinputlisting[language=R, firstline=179, lastline=180]{PS04_NF.R} 
	
	We get the following result:
	
	\begin{verbatim}
	> df2
	[1] 130
	\end{verbatim}
	
	We use these values to conduct a two-tailed t-test:
	\lstinputlisting[language=R, firstline=182, lastline=182]{PS04_NF.R} 
	
	We get the following result:
	
	\begin{verbatim}
	> 2*pt(t_stat2, df2, lower.tail=FALSE)
	[1] 0.001563955
	\end{verbatim}
	
	0.001563955 \textless 0.05 \\ 
	We can reject the null hypothesis that our coefficient is 0, meaning there is an association between being next to precincts with the yard signs and the vote share for Ken Cuccinelli.
	
	
	\vspace{7cm}
	\item [(c)] Interpret the coefficient for the constant term substantively. \\
	In a control precinct, we estimate the average proportion of the vote that went to Ken Cuccinelli is 30,2%.
	\vspace{1cm}
	
	\item [(d)] Evaluate the model fit for this regression.  What does this tell us about the importance of yard signs versus other factors that are not modeled? \\
	R² = 0.094 \\
	This is a very low result, showing that the model only explains 9.4\% of the total variation in vote share for Ken Cuccinelli. Either we need a better theoretical model to explain more or the variation, or we are satisfied with this result because we specifically wanted to look at the effect of the yard signs (which is the case in this study). This means that a large proportion of the voteshare can be explained by other factors (which makes sense). We can think of party affiliation, socio-economic status, gender, ideology, perception of the state of the economy...
	
	
	

	
	
	
\end{enumerate}  


\end{document}
